\newpage
\section{Podsumowanie}
Celem pracy było zbadanie, jak wybrane architektury głębokiego uczenia różnią się pod względem skuteczności oraz odporności na redukcję liczby danych treningowych i degradację ich jakości w zadaniu klasyfikacji chorób roślin. Dodatkowym celem było stworzenie prototypu, w postaci aplikacji mobilnej, która jest przykładem wykorzystania badanego zagadnienia.

W ramach pracy przeprowadzono przegląd literatury,dokonano wyboru zbiór danych oraz przeglądu różnych technik uczenia maszynowego oraz przedstawiono wybrane architektury poddanych analizie - ViT, EfficientNet, MobileNet,  ResNet, oraz DenseNet. Przeprowadzono szereg eksperymentów pozwalających ocenić wpływ jakości oraz ilości danych na skuteczność badanych modeli. Uzyskane wyniki zostały wykorzystane w ramach analizy porównawczej architektur oraz umożliwiły określenie różnic w ich odporności na ograniczenia związane z ilością i jakością danych treningowych.  

 Na podstawie przeprowadzonych badań można stwierdzić, że większość analizowanych architektur biorących udział w badaniu jest skuteczna w ramach zadania klasyfikacji chorób roślin. Jedynym wyjątkiem jest  rodzina modeli DenseNet, gdzie oba warianty osiągnęły znacznie gorsze wyniki od pozostałych badanych modeli w obu przeprowadzonych eksperymentach. 
 
 Na podstawie pierwszego eksperymentu ustalono, największą odporność na zmniejszanie liczby danych treningowe mają modele MobileNetV2, EfficientNetB0 oraz EfficientNetB1. Miały one najmniejszą różnicę miary Macro F1. Najlepszy wynik dla całego zbioru uzyskały modele ResNet, jednak miały one większe spadki jakości, niż wcześniej wspomniane modele. Najlepszym kompromisem między jakością klasyfikacji, odpornością na spadek liczby danych oraz wymaganiami pamięciowymi stanowiły modele z rodziny MobileNet, które osiągały zbliżone wyniki do modeli z większą ilością parametrów, a przy tym nie odbiegały jakością od innych analizowanych modeli.
 
 Na podstawie drugiego eksperymentu ustalono, że największą odporność na zniekształcone dane treningowe mają modele ResNet50 oraz EfficientNetB1. Różnica między Macro F1 dla użytych zbiorów treningowych była dla nich najmniejsza i wynosiła zaledwie kilka tysięcznych. Najlepszym kompromisem, między jakością klasyfikacji, odpornością na spadek jakości danych oraz rozmiarem były modele z rodzin EfficientNet i MobileNet - ich wyniki były zbliżone do innych z zestawienia, a jednocześnie miały one zdecydowanie mniej parametrów.
 
 W ramach finalnej analizy obu eksperymentów ustalono, że najlepszym kompromisem między wynikami obu eksperymentów, a wydajnością pamięciową są rodziny EfficientNet i MobileNet. Najlepsze wyniki osiągały modele z rodziny ResNet, a najgorzej z oboma zadaniami radziły sobie modele z rodziny DenseNet. Nie zaobserwowano przewagi modeli opartych na transformerach nad modelami CNN - modele z rodziny ViT osiągały zbliżone, choć gorsze wyniki od modeli CNN. W ramach badania nie zaobserwowano także zależnosci między rozmiarem modelu a wynikami. Największe modele, z rodziny ViT, nie uzyskały przewagi nad modelami mniejszymi, a uzyskały wyniki gorsze. 
 
 
 \subsection{Dalsze kierunki badań}
 Otrzymane rezultaty badań stanowią podstawę do dalszych badań nad doborem i optymalizacją architektur przedstawionych w badaniu. Poprzez zastosowanie identycznych parametrów podczas trenowania modeli uzyskano wyniki, które można ze sobą obiektywnie porównać, wybrane hiperparametry mogą nie być optymalne dla danych modeli. Interesującym kierunkiem byłoby sprawdzenie, czy odpowiednie dobranie hiperparametrów i zastosowanie dodatkowych technik augmentacji danych pozwoliłoby uzyskać przewagę modelowi ViT nad CNN. W przyszłych badaniach warto także zastosować inny zbiór danych - dane wykorzystane w tym badaniu są zebrane w jednym, kontrolowanym środowisku laboratoryjnym. Sprawdzenie, czy modele radzą sobie w detekcji chorób \newline w środowisku naturalnym, gdzie zdjęcia mogą być wykonane pod różnym kątem, z różnym tłem, czy obiektami zasłaniającymi część liści pozwoliłoby ocenić, czy  dane modele są \newline w stanie znaleźć zastosowanie w warunkach rzeczywistych, bezpośrednio w gospodarstwie rolnym, a nie w kontrolowanym środowisku laboratoryjnym.  Dalsze badania mogłyby także rozszerzyć zakres badanych modeli - badanie to obejmowało tylko 10 modeli, \newline a analiza większej liczby architektur, w tym tych opartych o transformery, pozwoliłaby na szerszą ocenę oraz sprawdzenie, czy przewaga modeli CNN nad ViT utrzymuje się również w przypadku innych rozwiązań.
  
 


